\documentclass[11pt]{article}
\usepackage[margin=1in]{geometry}
\usepackage{amsmath,amssymb,amsthm,booktabs,hyperref,xurl}
\setlength{\emergencystretch}{2em}
\newtheorem{theorem}{Theorem}\newtheorem{lemma}{Lemma}\newtheorem{corollary}{Corollary}
\newcommand{\F}{\mathbb F}\newcommand{\rank}{\operatorname{rank}}
\newcommand{\Hom}{\operatorname{Hom}}\newcommand{\GL}{\operatorname{GL}}
\newcommand{\gb}[2]{{#1\brack #2}_q}
\title{Acquired crossing phase carriers: exact counting through chains of diamonds}
\author{Astra P3 constructive counting research}\date{27 October 2026}
\begin{document}\maketitle
\begin{abstract}
The acquired matrix-dual counter previously summed arbitrary affine phases
when their phase-kernel image submodules were nested. This note removes that
condition completely for two matrix factors. More generally, it admits an
unbounded number of acquired crossing pairs, ordered by containment. Actual
binary linear algebra recognizes the required arrangement and manufactures
all intersections and sums. A finite-field graph count gives the transfer
across a crossing pair; its states are only an intersection dimension and a
rank polynomial. No subspaces or sequences of intersection choices are
enumerated. The resulting ordinary-circuit compiler is polynomial on the
certified family, including acquisition and exact reconstruction. It does
not give polynomial counting for arbitrary circuits or arrangements.
\end{abstract}

\section{The acquired data and the local signed polynomial}
Retain the actual algebra certificate of the matrix-factor construction.
The left nucleus is $S\simeq M_p(\F_q)$, $q=2^a$; the actual modules $U,O$
have binary dimension $apd$. The acquired reverse dual has dimension $c$ and
\[
 \mathcal D=LD_0,\qquad
 L=\bigoplus_{i=1}^bM_{s_i}(\F_q),\qquad
 d=\sum_i h_i s_i,\quad c=a\sum_i s_i^2.
\]
Here $D_0:O\to U$ is an isomorphism certified by regular-module dimensions
and the common-kernel test. It is used only in the proof; no cyclic element,
matrix units or field coordinates are searched for or supplied.

For the original target $y\in O$, each nonzero factor must occupy exactly
one joint $S,L_i$ simple-module copy. The actual row-space carrier
\[
 E_i=\sum_{D\in\text{an acquired basis of }\mathcal D_i}
       \operatorname{rowspace}_2(s\mapsto sD(y))\subseteq S^*
\]
is manufactured. Require that all nonzero $E_i$ coincide, with common
natural dimension $k=\dim_2E_i/(ap)$; write this right $S$-module as $R$.
Zero-target factors are kept separately.

The affine Hom offset gives a binary phase $\ell_i:\mathcal D_i\to\F_2$.
Manufacture the translates $D\mapsto\ell_i(AD)$ for an actual basis
$A\in L_i$. Their span has dimension $as_it_i$, where $t_i$ is the phase
matrix rank. Solve for their common annihilator $K_i\subseteq\mathcal D_i$,
and manufacture
\[
 F_i=\sum_{D\in\text{a basis of }K_i}
       \operatorname{rowspace}_2(s\mapsto sD(y))\subseteq R,
 \qquad f_i=\dim_2F_i/(ap).
\]
The shared acquisition checks the actual inclusion and the integral ranks;
it does not treat arbitrary binary spaces of divisible dimensions as
$S$-submodules. In proof coordinates $F_i=C_i\ker B_i^T$, for a target
matrix $C_i$ and phase matrix $B_i$. Thus
\[
 \max(0,k-t_i)\le f_i\le\min(k,s_i-t_i).
\]

For a candidate natural $j$-submodule $J\le R$, put
$\eta_i=\dim(J\cap F_i)$. The maps with target image contained in $J$ form
a rectangular matrix space of source dimension $s_i$ and target dimension
\[
 r_i=s_i-k+j,
 \qquad \tau_i=t_i-k+j+f_i-\eta_i
\]
for the rank of its restricted affine character. Let
$P_{s,r,\tau,h}(x)$ be its signed rank polynomial, with exponent
$h\rank A$. It is obtained without enumerating matrices from
\begin{equation}\label{eq:local-moments}
 m_v=\gb{r-\tau}{v}q^{sv},\qquad
 m_v=\sum_{u=0}^{v}P_u\gb{r-u}{v-u},\qquad
 P_{s,r,\tau,h}(x)=\sum_uP_ux^{hu}.
\end{equation}
These are image-containment moments: the character sum vanishes unless
the selected $v$-image space lies in its phase annihilator, and otherwise
equals $q^{sv}$. Lower triangular inversion proves the formula.
The affine-flag proof establishes the acquired identification of $F_i$ and
the displayed restricted phase rank; this extension uses the same code.

\section{Recognizing and manufacturing the chain}
Group all identical nonzero-target phase carriers $F_i$ by their canonical
actual binary row spaces. Add $0$ and $R$ as endpoints. For every pair test
both actual inclusions. Accept precisely when each distinct carrier is
incomparable with at most one other carrier. Repeated factors at one carrier
do not create extra vertices. This is a check on the acquired arrangement,
not an input decomposition or an assumed constant number of crossings.

\begin{lemma}[Ordered crossing pairs]\label{lem:chain}
Under this condition, the singleton carriers and incomparable pairs form
an ordered sequence. For every incomparable pair $F,G$, manufacture
$I=F\cap G$ and $T=F+G$. The unmatched carriers together with all such
intersections and sums form a chain. The endpoints $I,T$ of each pair are
adjacent in that chain.
\end{lemma}
\begin{proof}
A third represented carrier $H$ must be comparable with both $F$ and $G$.
If its comparison directions differed, transitivity would make $F,G$
comparable. Hence either $H\subseteq F\cap G$ or $H\supseteq F+G$.
Apply the same observation to a second incomparable pair. Its members
cannot be on opposite sides of the first pair, since they would then be
comparable. Therefore entire pairs are ordered. Inserting the intersection
and sum of each pair preserves this order. No other carrier or pair endpoint
can lie strictly between those two inserted endpoints. Binary intersections
and sums are right $S$-submodules because the original spaces are.
\end{proof}
The implementation also verifies every consecutive spine inclusion and
every claimed endpoint adjacency. These are polynomial checks and retain
explicit rejection diagnostics.

\section{The crossing-pair transfer}
All dimensions in this section are natural dimensions over $\F_q$; Morita
equivalence identifies the actual right $S$-submodule lattice with this
finite-vector-space lattice for the proof. No such coordinate identification
is computed by the algorithm.

Fix an incomparable pair $F,G$, set $I=F\cap G$, $T=F+G$, and write
\[
 i=\dim I,\quad A=\dim F-i,\quad B=\dim G-i,
 \quad\dim T=i+A+B.
\]
Suppose the current intersection $E=J\cap I$ is a fixed $e$-subspace. For
$J_T=J\cap T$ let
\[
 t=\dim J_T,\quad
 \alpha=\dim(J_T\cap F)-e,\quad
 \beta=\dim(J_T\cap G)-e,\quad
 u=t-e-\alpha-\beta.
\]
Use Gaussian coefficients equal to zero outside their usual ranges.

\begin{lemma}[Exact diamond multiplicity]\label{lem:transfer}
The number of choices for $J_T$, conditional on this fixed $E$, is
\begin{equation}\label{eq:transfer}
 q^{(i-e)(t-e)}
 \gb{A}{\alpha}\gb{B}{\beta}
 \gb{A-\alpha}{u}\gb{B-\beta}{u}
 |\GL_u(\F_q)|,
 \qquad |\GL_u(\F_q)|=\prod_{v=0}^{u-1}(q^u-q^v).
\end{equation}
This depends on the old intersection only through $e$.
\end{lemma}
\begin{proof}
Choose complements in the proof so that $T/I$ is a direct sum of spaces of
dimensions $A,B$. The projection $\overline{J_T}$ has dimension $t-e$;
its intersections with the two summands have dimensions $\alpha,\beta$.
Choose those two intersections in the first two Gaussian factors. After
quotienting by them, the remaining $u$-space meets both summands trivially.
It is uniquely a graph of an isomorphism between a $u$-space in each
remaining summand. These choices give the next two Gaussian factors and
$|\GL_u|$. Finally, a lift with fixed intersection $E$ is the graph of a
linear map $\overline{J_T}\to I/E$ after choosing a complement; there are
$q^{(i-e)(t-e)}$ such maps. The construction and its inverse are bijective.
\end{proof}
This is a classical finite-subspace graph count. The present claim is its
use after actual phase acquisition, not historical priority for Gaussian
coefficients or graph parametrizations. A closely related direct-sum
kernel census also occurs in the retained two-eigenspace quadratic source.

For comparison, an ordinary chain step from dimension $i$ to dimension $s$
has conditional multiplicity
\begin{equation}\label{eq:ordinary}
 \gb{s-i}{t-e}q^{(i-e)(t-e)}.
\end{equation}
The proof again chooses the quotient image and its graph into $I/E$.

\section{Dimension transfer and exact reconstruction}
Fix the final dimension $j$. Start at the zero spine node with intersection
dimension zero, and multiply the local polynomial for every factor whose
phase carrier is zero. Zero-target factors contribute their full signed
matrix-rank polynomials independently. At an ordinary spine step use
\eqref{eq:ordinary}. At a diamond step sum \eqref{eq:transfer} over
$e,t,\alpha,\beta$, multiplying the grouped factor polynomials at $F$ and
$G$ with intersection dimensions $e+\alpha,e+\beta$ respectively.
At every new spine node also multiply the factors carried by that node,
using intersection dimension $t$. Endpoints that are actual carriers are
therefore counted once, as are the two interior carriers of each diamond.

The state is a current intersection dimension and a polynomial in $x$.
The coefficient of $x^v$ records total represented dual rank $v$. States
with $t>j$ or $j-t>k-\dim T$ are discarded. Extension counts depend only on
the current intersection dimension, by Lemma~\ref{lem:transfer}, even if
earlier factors have weighted different actual intersections unequally.
Thus the transfer sums all histories without retaining or enumerating
their sequences. At the final node $R$, the intersection dimension is $j$.
Its polynomial is
\[
 M_j(x)=\sum_{\substack{J\le R\\\dim J=j}}
       \sum_{D:\,\operatorname{carrier}(Dy)\subseteq J}
            (-1)^{\ell(D)}x^{\rank_2D/(ap)}.
\]

\begin{theorem}[Acquired chain-of-diamonds counter]
Let $H_\rho(x)$ be the polynomial for dual elements whose target carrier
has exact dimension $\rho$. It is recovered by the triangular inversion
\[
 H_j(x)=M_j(x)-\sum_{\rho<j}\gb{k-\rho}{j-\rho}H_\rho(x).
\]
If $H_{v,\rho}$ denotes its coefficients, the exact original-input count is
\begin{equation}\label{eq:answer}
 \boxed{2^{-c}\sum_{v,\rho}H_{v,\rho}
        F_{p,d,d,q}(k,v,\rho).}
\end{equation}
Here $F$ is the proved signed universal-product fiber. All algebra,
phase, arrangement and reconstruction work is deterministic polynomial
time in the explicit ordinary-circuit description length.
\end{theorem}
\begin{proof}
A fixed $\rho$-carrier lies in exactly
$\gb{k-\rho}{j-\rho}$ of the $j$-submodules. The transfer has counted every
actual intermediate intersection precisely once, so double-counting gives
$M_j=\sum_{\rho\le j}\gb{k-\rho}{j-\rho}H_\rho$.
Triangular inversion has unit diagonal and is exact over the integers.
The acquired Fourier formula groups characters by precisely represented
rank and target rank, giving \eqref{eq:answer}. The frontend's affine
input reconstruction is bijective; no internal program variables are
counted as independent inputs. The implementation checks the signed total
mass, the zero character and final divisibility by $2^c$.
\end{proof}

\section{Costs and ordinary-input families}
Let $L$ be the explicit ordinary-DAG description length and
$N=n+m+G+2=O(L)$, where $n$ is the declared number of independent inputs,
$m$ the number of marked outputs and $G$ the gate count. All natural
dimensions and numbers of factors are $O(N)$; actual algebra dimensions
are $O(N^2)$. There are $O(N)$ acquired carrier nodes and diamonds.
Acquisition uses the existing polynomial actual-algebra certificates;
pairwise inclusions, intersections and sums use binary linear algebra.

For each $j$ the transfer keeps at most $k+1$ rank polynomials of degree
at most $d$. A diamond uses at most $O(N^4)$ dimension-index transitions.
With $O(N)$ spine steps and $O(N)$ final dimensions, this is $O(N^6)$
transitions, independent of the number of subspaces or histories.
Gaussian coefficients, signed local sums, transfer intermediates and final
fiber terms have $O(N^2)$ bits. Naive polynomial multiplication and integer
arithmetic, together with acquisition, are covered by the conservative
uniform bound $O(N^{20})$ bit operations and polynomial space. The final
answer uses at most $n+1$ bits. This bound makes no claim of practical
superiority over equally informed specialized competitors.

\begin{corollary}[Two factors: every affine phase]
For exactly two matrix factors with a common nonzero target carrier and
the certified one-copy targets, every affine Hom offset is admitted.
\end{corollary}
\begin{proof}
There are at most two distinct proper phase carriers. They are either
nested or one incomparable pair; both configurations pass the test.
\end{proof}

A concrete all-size ordinary-input family uses $q=2^a$, $s\ge2$, $p=s$,
$d=4s$, and
\[
 \mathcal D=\{\operatorname{diag}(A,A,B,B):A,B\in M_s(\F_q)\},
 \quad W=\mathcal D^\perp.
\]
Right inputs range bijectively over $Y_0+W$ for an arbitrary fixed matrix
$Y_0$; left inputs range over $X\in M_{s\times4s}(\F_q)$. The output is
$X(Y_0+Y)$, with target $(C_0,0,C_0,0)$ for arbitrary $C_0$.
The original independent-input count and codimension are
\[
 n=18as^2,\qquad c=2as^2=n/9.
\]
The raw-product-family proof acquires the polar split, full matrix nucleus
and actual repeated-block dual. Adding $Y_0$ changes only the linear term;
it is an actual Hom offset and is handled by the affine acquisition.
Both target carriers coincide. Thus arbitrary affine phases are handled
for every size, including phases whose two carriers cross. Dense binary
input conjugacy, output mixing and the recognized reversible suffix retain
the count; admission is limited to the specified exact quadratic
propagation and terminal syntax, not arbitrary equivalent encodings.

For unbounded numbers of diamonds take $b=2r$, $s\ge2r$, the repeated-block
family with $d=2bs$, and common full-rank target. Choose phase kernels
\[
 \langle e_1,\ldots,e_{2u},e_{2u+1}\rangle,
 \quad\langle e_1,\ldots,e_{2u},e_{2u+2}\rangle,
 \qquad0\le u<r.
\]
These are realized by diagonal phase matrices and hence by ordinary affine
right offsets. The number of diamonds, matrix dimension, codimension and
original input dimension all grow. The actual compiler manufactures this
structure after hiding; no flag or diamond decomposition is supplied.

\section{Executed checks and explicit limits}
\texttt{phase\_diamond\_audit.py} performs literal finite-field subspace
censuses, literal matrix/image-space convolution, and ordinary-gate tests.
The saved JSON distinguishes native acquisition from formula-only growth.
The first completed checks cover 231 canonical pair types over
$\F_2,\F_4,\F_8$, comparing 9,183 exact census coefficients.
A two-diamond calculation on four $4\times4$ binary factors matches
independent enumeration of all 65,536 local matrices and convolution of
their actual image subspaces; it has 35 nonzero rank profiles.
A three-diamond, six-factor formula execution has parameter $n=5400$,
$c=216$ and a 4,968-bit answer. No raw 5,400-input acquisition is claimed.

All 256 actual affine phase words of the two-factor $2\times2$ native
example agree with literal enumeration of all 256 acquired dual elements
and the independent prior character sum. This includes all 54 phases
declined by the nested-carrier counter. A newly generated 72-input circuit
with 24,676 ordinary gates and 500 nonlinear reversible updates acquires
two crossing phase lines. Its exact count is
\[
 71{,}210{,}970{,}084{,}540{,}416.
\]
An independent oracle enumerates all 65,536 original left matrices and
solves the actual affine right-input equations, obtaining the same count.
Another 128 original-gate evaluations match the explicit planted map.
The final wrapper using the shared actual-phase acquisition was replayed
against the native certificate and agrees exactly. These are correctness
checks, not a runtime comparison against optimized competitors.

Three incomparable lines are retained
as a decisive rejected configuration: a carrier has two incomparable
neighbors, so this transfer is inapplicable. More general phase arrangements
and distinct nonzero target carriers are not silently counted by this
theorem. In particular this work does not establish polynomial exact
counting for arbitrary Boolean circuits.

\section*{Implementation and sources}
\begin{itemize}
\item \texttt{work\_phase\_diamond/phase\_diamond.py}: actual arrangement
recognition, dimension transfer and native-count wrapper.
\item \texttt{work\_matrix\_root/matrix\_affine\_flag\_counter.py}:
shared actual phase acquisition and local rectangular signed moments.
\item Retained proofs: \emph{Acquired matrix-factor actions and interacting
target geometry}; the affine-flag theorem; and
\emph{Raw product family for acquired semisimple matrix duals}.
\item H. Cohn, \emph{Projective geometry over $\F_1$ and the Gaussian
binomial coefficients} (2004), for the classical finite-subspace meaning
of Gaussian coefficients: \url{https://arxiv.org/abs/math/0407093}.
\end{itemize}
\end{document}
